\section{Resultados e Discussão}

Os resultados obtidos durante os dois primeiros trimestres de execução do
projeto encontram-se integralmente documentados em repositório público
versionado, desenvolvido com o objetivo de garantir transparência,
rastreabilidade e reprodutibilidade científica das etapas computacionais
realizadas ao longo da iniciação científica \cite{leao2026vqe}.

O repositório reúne notebooks experimentais, códigos modulares em Python,
registros das simulações quânticas e dados utilizados nas etapas de
validação clássica, constituindo o ambiente central de desenvolvimento
do presente trabalho.